% RESEARCH_DRAFT_SHA256: 41ca9ae3c50ff5886b7ba60cd659163b689f8d39ff3d8718e975e871ab051c00
% Layout and source-reference wrappers expanded for the web renderer.
\documentclass[11pt]{article}
\usepackage[margin=25mm]{geometry}
\usepackage{amsthm}
\usepackage{fontspec,unicode-math}
\setmainfont{Latin Modern Roman}
\setmathfont{Latin Modern Math}
\usepackage{amsmath,amssymb,mathtools,xurl,hyperref}
\hypersetup{colorlinks=true,urlcolor=blue,linkcolor=blue}
\setcounter{secnumdepth}{0}
\setlength{\parindent}{0pt}
\setlength{\parskip}{5pt}
\setlength{\emergencystretch}{3em}
\newcommand{\sref}[2]{\hyperlink{source:#1:#2}{[S#2]}}
\newcommand{\cataloguescope}[1]{\paragraph{Source record.}#1}
\newcommand{\floor}[1]{\left\lfloor#1\right\rfloor}
\newtheorem{proposition}{Proposition}
\newtheorem{lemma}[proposition]{Lemma}
\begin{document}
\paragraph{Research manuscript.} OpenAI \texttt{gpt-6-astra}, reasoning \texttt{ultra}; 11 October 2026.
\paragraph{Preserved source body.} The complete catalogue mathematical text and sources follow with layout and source links expanded, including source qualifications and its historical preparation notes. New research begins at the section headed Research attempt.

% UnsolvedMath ID 20002123; source catalogue code AIM-INFRASTRUCTURE-0023
\section{Separating exposure from retained mastery under a curriculum time budget}\label{20002123}
{\small UnsolvedMath ID 20002123 · Computer Science}
\subsection{Definitions and mathematical statement}

\paragraph{Shared notation.}Standard number systems are $\mathbb N,\mathbb Z,\mathbb Q,\mathbb R,\mathbb C$. Specialized notation and index ranges follow the cited source and the definitions below.


Active learning asks students to carry out mathematical activities during instruction. Curriculum coverage measures which planned topics receive instructional time; demonstrated learning measures what students can subsequently do or explain. Topic coverage, depth of understanding, and time spent are separate quantities.
\paragraph{Record scope and qualifications.}
The catalogue title is retained as metadata. Named certificates, contracts, bounds or frameworks in that title are not asserted as results of the original source extract.
The registry identifies reconstructed wording. The premise about a mistaken impression is preserved as an educational claim requiring evidence, not asserted here as an established result.

Catalogue screening limitation: Educational assessment, adoption, or study-design discussion; the supplied question does not state a mathematical target with defined assumptions and outcome.

The source snapshot labels this record partially solved; this is catalogue metadata rather than independent verification of a mathematical result or the current completion of a project.
The source text below is rendered literally where necessary; damaged formulas and composite extracts are not silently repaired into a new mathematical claim.
The supplied catalogue marks this formulation as an unverified reconstruction; missing assumptions are not certified by the source link.
\paragraph{Statement recorded in the supplied source.}
25. The mistaken impression that active learning does not cover enough material

(counteract the "curriculum focus") \href{http://aimath.org/pastworkshops/oerassessmentproblems.pdf}{[S2]}
\paragraph{Source status.}
The supplied catalogue labels this source record partially solved; that metadata was checked against the supplied snapshot on 2026-10-11. The label is preserved for provenance and is not an independent verification of present open status or mathematical correctness.
\subsection{Short English statement}
Examine and address the perception that active learning covers insufficient material.
\subsection{Sources}
\begin{itemize}
\item[S1] Ulam AI and the UnsolvedMath Contributors, UnsolvedMath: A Curated Collection of Open Mathematics Problems, record 20002123 (AIM-INFRASTRUCTURE-0023). Catalogue curation and metadata: CC BY 4.0. The source extract is retained; definitions, English summaries and qualifications are editorial additions.
\url{https://huggingface.co/datasets/ulamai/UnsolvedMath}.
\item[S2] AIM Workshop Problem Lists, source linked by catalogue record AIM-INFRASTRUCTURE-0023.
\url{http://aimath.org/pastworkshops/oerassessmentproblems.pdf}.
\end{itemize}
\paragraph{Preparation and review.}
The original source and UnsolvedMath contributors retain attribution for the source extract. Codex (GPT-6; exact serving revision not exposed) assisted with definitions, English summaries, scope notes and formatting. No new human mathematical authorship or independently verified solution is asserted. Permanent source ID 20002123 and code AIM-INFRASTRUCTURE-0023 checked against the complete supplied registry and current repository. Distinct upstream records, including repeated questions and broad or damaged extracts, are retained for complete coverage of the source CS category. A repeated record is not a new independent conjecture. Complete mathematical-equivalence and current-literature review has not been performed. Preparation and final review follow contribution prompts C and D, with source gaps disclosed. Structural and rendering checks do not establish a new theorem or open status.
\label{end:20002123}

\section{Research attempt}
\subsection{Literature, source context and scope}
The primary AIM notes describe a brainstorming session. The original excerpt and all catalogue qualifications remain above. They do not specify a mathematical theorem or a complete empirical study. The formal model below is an explicitly limited research route introduced for this attempt, not a reconstruction of a missing original conjecture. The supplied catalogue's earlier research summary was inspected as context; its named framework is not adopted as an established solution. The elementary result below is not claimed novel, and no educational effectiveness data were collected.

\subsection{Attempt: does equal topic coverage imply equal retained coverage?}
The source's assertion about a mistaken impression is retained without treating it as established evidence. I tested a limited scheduling model: $n$ lessons have equal duration, complete at fixed increasing times $t_1<\cdots<t_n\le T$, and topic $i$ has weight $w_i\ge0$. Its retention probability at time $T$ is stipulated to be $e^{-\lambda(T-t)}$ when taught at time $t$, with common $\lambda>0$. All schedules cover exactly the same topics.
\begin{proposition}
The weighted expected retained total is maximized by placing topic weights in nondecreasing order of completion time. Equal exposure counts alone do not fix this retained total.
\end{proposition}
\begin{proof}
Let $r_j=e^{-\lambda(T-t_j)}$, so $r_j$ is strictly increasing. If an earlier position $a$ has weight $u$ and a later position $b$ has smaller weight $v$, their current contribution is $ur_a+vr_b$. Swapping them changes it by
\[
(vr_a+ur_b)-(ur_a+vr_b)=(u-v)(r_b-r_a)>0.
\]
Removing inversions repeatedly yields nondecreasing weights without decreasing the objective; a finite sequence of adjacent swaps suffices. Every order without inversions therefore attains the maximum, including arbitrary permutations of equal weights. Since the displayed difference can be strictly positive, two schedules with identical exposure counts can have different expected retained totals.
\end{proof}
\subsection{Verification and remaining gaps}
For weights 1 and 3, completion times 1 and 2, $T=2$ and $\lambda=\log2$, the ordered total is $1/2+3=7/2$, while the reversed total is $3/2+1=5/2$. Normalizing by total weight four gives $7/8$ versus $5/8$. The failed inference is that mentioning the same topics settles a retention comparison. The exponential law, equal durations and weights are assumptions; prerequisites or different forgetting rates can change the scheduling problem. This supplies no empirical active-learning superiority result and no answer about actual curriculum time tradeoffs.

\section{References and source attribution}
American Institute of Mathematics, workshop discussion notes on open educational resources and assessment, \url{https://aimath.org/pastworkshops/oerassessmentproblems.pdf}. The 16-page primary PDF was opened on 11 October 2026 and its numbered brainstorming item was checked against the supplied record. The source supports the provenance and agenda, not the additional modeling assumptions or an empirical educational conclusion. Catalogue and contributor attribution are retained in the reproduced Sources section. The worked mathematics is proved here without reliance on an unverified external theorem.


\paragraph{Integration review.} Codex (GPT-6; exact serving revision not exposed) reviewed the written arguments and expanded layout and source-reference wrappers. Repository metadata and the production web renderer were checked using Python 3.13 and JavaScript with MathJax 3.2.2. This model-assisted review is not independent human proof verification; no PDF compilation or formal proof checking is asserted.

\section{Completion Estimate}
\noindent\textbf{COMPLETION: 0\%}
No substantiated new progress toward the entire original target is claimed. The worked results and failed route are delimited above; a known result, illustrative model, or auxiliary obstruction is not counted as original-target completion.

\section{Final status and attribution}
\textbf{Attempt status: unresolved.} This describes this research attempt, not a new ruling on the source catalogue's status or on external literature claims.
The mathematics and this self-audit were generated by OpenAI \texttt{gpt-6-astra}, reasoning \texttt{ultra}, on 11 October 2026. Codex (GPT-6; exact serving revision not exposed) orchestrated the batch and contributed formatting/repository review. No human mathematical authorship, independent human proof checking, or formal verification is claimed. The historical preparation paragraphs in the reproduced catalogue body concern that body only.

Contribution prompt A was applied to the complete definition and supplied source record, followed by the adversarial checks of prompt B. Local reproduction files are \texttt{batch008/checks.py}, \texttt{computations.json}, and the per-ID audit JSON. They were not public links at preparation time. Numerical checks are finite and do not prove unrestricted results. No LaTeX engine was found locally; compilation was not performed by this research agent.

\end{document}
